% !TEX root = ../notes.tex

\chapter{符号定义}

% \epigraph{数学是描述世界的语言，但就像世界上大多数语言一样，总有人能用一些合理的符号组成令人迷惑的句子.}{——ZexHexoft\\于2026.9.8晚上23:22}

书中主要采用与David J.Griffiths所著的《电动力学导论(Introduction To Electrodynamics)》类似的记号(但有细微的不同).

具体而言：
\begin{itemize}
    \item 在这里，我们使用诸如$\mathbf{r}$,$\mathbf{v}$,$\mathbf{a}$的符号表示矢量
    \item 而$r,v,a$则对应他们的"大小"
    \item $\hat{\mathbf{r}}$,  $\hat{\mathbf{v}}$,  $\hat{\mathbf{a}}$表示同方向的单位向量
    \item 特殊的，我们一般用$\mathbf{r}$表示位置矢量(由原点指向该位置), 用$\mathbf{v}$表示速度
    \item 更特殊的，我们使用$\brcurs$ 表示"间隔矢量"(即$\brcurs=\mathbf{r}-\mathbf{r}'$, 从源点指向场点), 用$\rcurs$表示其大小
    \item 对于某个量(例如$r,\mathbf{v}$)对时间的一阶导数，我们一般写成$\dot{r},\dot{\mathbf{v}}$的形式, 求几阶导就加几个点, 甚至可以$\overset{........}{y}$ 
    \item 在积分运算中，我们会使用,$\mathrm{d}\mathbf{a}$或者$\mathrm{d}\mathbf{S}$表示面积微元, 用$\mathrm{d}\tau$或者$\mathrm{d}^3\vec x$表示体积微元.
    \item 另外，我们有时会使用爱因斯坦求和规则，阅读时请注意.
\end{itemize}

有人可能会觉得这里的符号有点奇怪(也许确实是这样), 但谁让我读的第一本电动力学教材是Griffiths的, 而你又选择了读这本笔记呢?
\footnote{当然，这里的所有记号规范可能只是我的一厢情愿，实际请以文意为准.}


